% --- Perfil ---
\vspace{-15pt}
\section{Perfil}
\begin{onecolentry}
	Soy graduado en Ingeniería de Telecomunicación, habiendo cursado el grado y el máster en la Universidad de Alcalá. Dentro del campo de la ingeniería, estoy especialmente interesado en las redes, los servicios y la seguridad, habiendo orientado mi trayectoria profesional hacia la operación y documentación de infraestructuras de red, la administración de sistemas Linux empresariales y la gestión de mecanismos de autenticación y protección de servicios. A lo largo de mi experiencia, he trabajado con entornos corporativos en los que he participado en tareas de despliegue, integración y automatización, combinando conocimientos de redes, sistemas y prácticas DevOps para mejorar la fiabilidad, la trazabilidad y el mantenimiento de distintos servicios y plataformas.
\end{onecolentry}

\vspace{-10pt}

% --- Educación ---
\section{Educación}

\begin{onecolentry}
	\begin{tabularx}{\textwidth}{@{}>{\raggedright\arraybackslash}X>{\raggedright\arraybackslash}p{3.6cm}>{\raggedleft\arraybackslash}p{2.4cm}@{}}
		\textbf{Máster en Ingeniería en Tecnologías de Telecomunicación} & Universidad de Alcalá & \textit{Febrero 2026}
	\end{tabularx}
	
	\vspace{0.05 cm}
	\hspace{0.2cm}Premio extraordinario
\end{onecolentry}

\vspace{0.2 cm}

\begin{onecolentry}
	\begin{tabularx}{\textwidth}{@{}>{\raggedright\arraybackslash}X>{\raggedright\arraybackslash}p{3.6cm}>{\raggedleft\arraybackslash}p{2.4cm}@{}}
		\textbf{Grado en Ingeniería en Tecnologías de Telecomunicación} & Universidad de Alcalá & \textit{Julio 2024}
	\end{tabularx}
	
	\vspace{0.05 cm}
	\hspace{0.2cm}2.º puesto de promoción y premio extraordinario
\end{onecolentry}

\vspace{-7pt}

% --- Experiencia ---
\section{Experiencia}

\begin{twocolentry}{\textit{Septiembre 2024 – Enero 2026}}
	\textbf{Infraestructuras, automatización y servicios} \\
	Universidad de Alcalá, Alcalá de Henares
\end{twocolentry}
\begin{onecolentry}
	\begin{highlights}
		\item Diseñé y desplegué servicios en contenedores, utilizando Docker y containerd, así como orquestadores como Docker Compose, Kubernetes y OpenShift.
		\item Integré flujos de despliegue automatizados en entornos empresariales mediante GitHub Workflows y Vagrant, tanto para entornos de prueba como de producción.
		\item Desarrollé pipelines y utilidades en Python, MATLAB y C para dar soporte a servicios técnicos y cargas intensivas.
	\end{highlights}
\end{onecolentry}

\vspace{0.2 cm}

\begin{twocolentry}{\textit{Febrero 2024 – Agosto 2024}}
	\textbf{Redes y sistemas} \\
	Correos y Telégrafos S.A., Madrid
\end{twocolentry}
\begin{onecolentry}
	\begin{highlights}
		\item Trabajé en la operación y documentación de equipos y servicios de red en entorno corporativo.
		\item Participé en tareas asociadas a routing, switching y administración de infraestructuras basadas en tecnologías Cisco y Fortinet.
		\item Gestioné servicios de seguridad y autenticación como RADIUS, LDAP, OAuth2 y TLS.
		\item Participé en despliegues y prácticas DevOps, integrando automatización en pipelines de CI/CD.
		\item Administré sistemas Linux empresariales, principalmente RHEL y Amazon Linux 2.
	\end{highlights}
\end{onecolentry}

\vspace{0.2 cm}

\begin{twocolentry}{\textit{Enero 2023 – Enero 2024}}
	\textbf{Ingeniería e investigación aplicada} \\
	Universidad de Alcalá / Correos y Telégrafos S.A., Alcalá de Henares
\end{twocolentry}
\begin{onecolentry}
	\begin{highlights}
		\item Desarrollé plataformas web de análisis de datos, basadas en arquitectura MVC, para proyectos técnicos y de investigación.
		\item Implementé soluciones de computación multihilo y en clúster para tareas de procesamiento intensivo.
		\item Me encargué de la administración de entornos Linux y de sistemas de computación en red.
		\item Colaboré en la actividad científica del grupo y participé como coautor en publicaciones \href{https://orcid.org/0009-0000-6862-1872}{ORCID}.
	\end{highlights}
\end{onecolentry}

\vspace{-10pt}

% --- Aptitudes ---
\section{Aptitudes}

\begin{onecolentry}
	\textbf{Programación:} C/C++, Python, Bash, HTML/CSS/JS, Groovy
\end{onecolentry}
\begin{onecolentry}
	\textbf{Redes y servicios:} Linux (Ubuntu, RHEL, AL2), routing, switching, Cisco, Fortinet, GNS3, Kubernetes, Red Hat OCP, vSphere ESXi, AWS
\end{onecolentry}
\begin{onecolentry}
	\textbf{Seguridad y autenticación:} FortiGate, RADIUS, LDAP, OAuth2, TLS
\end{onecolentry}
\begin{onecolentry}
	\textbf{CI/CD:} Jenkins, GitHub/GitLab, Vagrant, Ansible, Terraform
\end{onecolentry}